% Chapter 3: User Research

\chapter{User Research}\doublespacing

\label{Chapter3}

\lhead{Chapter III. \emph{User Research}}

%----------------------------------------------------------------------------------------
%	SECTION 3.1
%----------------------------------------------------------------------------------------

\section{Overview Of The Research}

This chapter grounds the problem space established in Chapter~1 in primary data collected directly from the IITK undergraduate population. Where Chapter~1 argued --- from general career-guidance literature and a first read of student sentiment --- that the four-year IITK lifecycle produces a specific pattern of choice-overwhelm, unclear prioritization, and unmet support, this chapter tests that argument against a structured sample of actual respondents, and uses the resulting patterns to define the stakeholders, personas, and design brief that Chapter~4 builds the game around.

Data was collected through a structured, self-administered online questionnaire \textit{[confirm platform --- Google Forms / Tally / other]}, distributed \textit{[confirm channels --- e.g.\ hostel-wing WhatsApp groups, batch mailers, Instagram broadcast]} over \textit{[confirm collection window, e.g.\ ``a two-week period in Month, Year'']}. Participation was voluntary and unincentivized \textit{[confirm if any incentive was offered]}. The questionnaire had four parts: (a)~demographics --- Year, Branch, and Gender; (b)~fifteen five-point Likert-scale statements (1~=~Strongly Disagree to 5~=~Strongly Agree) covering choice-overwhelm, time-balance, decision clarity, peer influence, existing support structures, wellbeing cost, and appetite for a game-based intervention; (c)~three multiple-choice questions on time allocation and decision-making reliance; and (d)~a closing consent question asking whether respondents were willing to be contacted for future playtesting.

The form received 124 raw submissions. Ten were exact duplicates --- identical Year, Branch, Gender, and submission-timestamp combinations, consistent with accidental double-submits rather than genuine repeat respondents --- and were removed, leaving $N = 114$ unique responses used throughout this chapter.

This phase relied on quantitative self-report only. Two to four short follow-up interviews were considered, in line with the \textit{Enrich} thesis this project's structure is modelled on, but were not conducted within this phase; this is noted as a limitation in Chapter~6 and as a natural next step, particularly since 58 of 113 respondents who answered the consent question (51.3\%) indicated willingness to be recontacted for playtesting or further research.

%----------------------------------------------------------------------------------------
%	SECTION 3.2
%----------------------------------------------------------------------------------------

\section{User Sample / Participants}

The 114 respondents were distributed unevenly across years (Figure~\ref{fig:chap3_fig3_1}): 90 First Years (79.0\%), 12 Alumni (10.5\%), 7 Fourth Years (6.1\%), and 5 Second Years (4.4\%), with no Third Year respondents. This skew toward First Years is a meaningful limitation on how far the findings below can be generalized across the full undergraduate lifecycle --- see the discussion in Section~3.4 and the caveat repeated in Chapter~6 --- but it is also, in one sense, on-topic: First Years are the population closest to the choice-overwhelm and orientation gap this project is centrally concerned with, and the 12 Alumni responses provide a small but valuable retrospective, hindsight-driven contrast used directly in the Section~3.6 personas.

\begin{figure}[htbp]
  \centering
  \includegraphics[width=\textwidth]{Images/chap3_fig3_1.png}
  \caption{Year-wise distribution of respondents ($N=114$).}
  \label{fig:chap3_fig3_1}
\end{figure}

%----------------------------------------------------------------------------------------
%	SECTION 3.3
%----------------------------------------------------------------------------------------

\section{User Research Findings}

\textit{[Note: This section contains the full survey analysis, including Likert-scale results for all 15 statements, multiple-choice distributions, and cross-tabulated findings. Due to the extensive nature of the base64-encoded figures in the source document, figure placeholders are used below. Replace each with the actual extracted image file.]}

\subsection{Likert-Scale Findings}

The fifteen Likert-scale statements were grouped into seven thematic clusters. The key findings are summarized below, with each cluster's agreement rates (sum of ``Agree'' and ``Strongly Agree'') and disagreement rates (sum of ``Disagree'' and ``Strongly Disagree'') reported.

\begin{itemize}
  \item \textbf{Choice-Overwhelm:} 82.5\% of respondents agreed or strongly agreed that IITK presents an overwhelming number of choices in academics, clubs, and career paths.
  \item \textbf{Time-Balance Difficulty:} 77.2\% agreed or strongly agreed that balancing academics, extracurriculars, and personal wellbeing is difficult.
  \item \textbf{Decision Clarity:} Only 28.1\% felt confident they understood the long-term consequences of their semester-by-semester choices.
  \item \textbf{Peer Influence:} 64.9\% agreed that peer decisions significantly influenced their own academic and career choices.
  \item \textbf{Existing Support:} Only 35.1\% agreed that existing institutional support (faculty advisors, CDC workshops) adequately helped them navigate the four-year lifecycle.
  \item \textbf{Wellbeing Cost:} 71.9\% agreed that the pressure of making consequential decisions without clear guidance negatively affected their mental health or stress levels.
  \item \textbf{Game-Based Intervention Appetite:} 78.9\% agreed or strongly agreed that a simulation game allowing them to ``rehearse'' the four-year IITK lifecycle would be a useful learning tool.
\end{itemize}

\begin{figure}[htbp]
  \centering
  \includegraphics[width=\textwidth]{Images/chap3_fig3_2.png}
  \caption{Likert-scale response distributions across all 15 statements.}
  \label{fig:chap3_fig3_2}
\end{figure}

\subsection{Multiple-Choice Findings}

The three multiple-choice questions addressed time allocation priorities, decision-making reliance, and retrospective regret.

\begin{figure}[htbp]
  \centering
  \includegraphics[width=\textwidth]{Images/chap3_fig3_3.png}
  \caption{Multiple-choice response distributions.}
  \label{fig:chap3_fig3_3}
\end{figure}

%----------------------------------------------------------------------------------------
%	SECTION 3.4
%----------------------------------------------------------------------------------------

\section{Sample Limitations}

The year-wise skew described in Section~3.2 is the most significant limitation of this dataset. With 79\% of respondents being First Years, the findings above are best understood as a snapshot of the population most immediately experiencing the orientation gap --- not as a representative cross-section of the full four-year student body.

A second limitation is the absence of qualitative data: the survey was entirely quantitative, and while two to four follow-up interviews were planned, they were not conducted within this research phase. This means the ``why'' behind many of the agreement rates remains unexplored; the numbers tell us \textit{that} students feel overwhelmed, but not the specific mechanisms or moments that produce that feeling. This is noted as a natural next step in Chapter~6.

%----------------------------------------------------------------------------------------
%	SECTION 3.5
%----------------------------------------------------------------------------------------

\section{Key Takeaways}

The survey data supports three central claims that directly inform the game design in Chapter~4:

\begin{enumerate}
  \item \textbf{Finding 1: Choice-overwhelm is real and widespread.} More than four in five respondents report feeling overwhelmed by the number of choices IITK presents.
  \item \textbf{Finding 2: Support structures are perceived as inadequate.} Nearly two-thirds of respondents do not feel that existing institutional support helps them navigate the lifecycle.
  \item \textbf{Finding 3: There is strong appetite for a game-based intervention.} Nearly four in five respondents expressed interest in a simulation game that lets them rehearse the four-year lifecycle.
\end{enumerate}

These three findings together constitute the primary evidence base for the design strategy statement in Section~3.7 and the game concept developed in Chapter~4.

%----------------------------------------------------------------------------------------
%	SECTION 3.6
%----------------------------------------------------------------------------------------

\section{User Personas}

Based on the survey data and the three key findings above, three composite personas were developed to represent the primary user archetypes the game is designed for.

\subsection{Persona 1: The Overwhelmed First Year}
\textit{[Based on the dominant survey cohort --- First Year students who report high choice-overwhelm and low decision clarity.]}

\subsection{Persona 2: The Retrospective Alumni}
\textit{[Based on the 12 Alumni respondents who bring hindsight-driven perspective on the four-year lifecycle.]}

\subsection{Persona 3: The Balanced Generalist}
\textit{[A composite persona representing students who attempt to balance academics, extracurriculars, and wellbeing --- the strategy the game is designed to make visible and testable.]}

%----------------------------------------------------------------------------------------
%	SECTION 3.7
%----------------------------------------------------------------------------------------

\section{Design Strategy Statement}

Drawing on the three key findings from Section~3.5 and the three personas from Section~3.6, the design strategy statement for this project is:

\begin{quote}
``How might we give IITK undergraduates a low-stakes way to rehearse the trade-offs of their four-year lifecycle --- balancing academics, clubs, projects, social life, and wellbeing against the ticking clock of eight semesters --- so that they can make more informed, less anxious decisions in real life?''
\end{quote}

This statement directly frames the game concept developed in Chapter~4.

%----------------------------------------------------------------------------------------
%	SECTION 3.8
%----------------------------------------------------------------------------------------

\section{Decision-Tree Analysis of IITK Placement Data}

This section describes a secondary data analysis conducted to ground the game's placement mechanics in empirical evidence. A decision-tree model was trained on 57 anonymized IITK resumes matched against their actual 2024 SPO placement outcomes. The five features that best separated one placement outcome from another were: CPI, Algorithmic/Coding skill, Positions of Responsibility (POR), Research experience, and Product/Design experience.

These five dimensions became the five resume-stat axes tracked in the game's internal economy (Section~4.6) and the basis for the placement eligibility thresholds in \texttt{placement.csv} (Section~5.1.7).

\begin{figure}[htbp]
  \centering
  \includegraphics[width=\textwidth]{Images/chap3_fig3_4.png}
  \caption{Decision-tree analysis: feature importance for placement outcome prediction.}
  \label{fig:chap3_fig3_4}
\end{figure}
